\section{Qualidade dos registros e cuidados adotados}
\label{sec:qualidade}

Nenhuma base administrativa é perfeita, e parte do trabalho
exploratório é localizar exatamente onde estão os problemas para que
eles não contaminem as conclusões. A boa notícia: os problemas deste
painel são poucos, identificados e circunscritos.

\textbf{Ocupação acima de 100\%.} Na internação, 15,5\% das observações
(492) ultrapassam 100\%, e na UTI, 6,5\% (207). Em sua maioria isso é
informação verdadeira de sobrecarga (mais diárias que a capacidade
cadastrada comporta) e não deve ser apagado. O problema real está na
cauda extrema da UTI: \textbf{os 15 maiores valores pertencem todos a
2020 e 2021}, liderados pelo \textbf{Hospital Irmãos Penteado e Santa
Casa de Campinas (CNES 2022648)}, que registra 875\% de ocupação de UTI
em 2021: são 6.389 diárias divididas pelos \textbf{2 leitos SUS de UTI}
então cadastrados, pois a razão do SIH é computada sobre os leitos SUS
(o cadastro total do hospital indicava 16 leitos de UTI, que ainda
assim dariam 109\% de ocupação). A
explicação é operacional: na onda da pandemia os hospitais expandiram
leitos de UTI às pressas e o cadastro não acompanhou, produzindo
divisões por um denominador defasado. Parte dos valores acima de 100\%
tem origem em um \textbf{descasamento de calendário}: leitos que
entraram no cadastro apenas no meio do ano (caso típico das UTIs
expandidas na pandemia) convivem com diárias acumuladas em período
maior que o de cadastro, inflando a razão. Esse viés não é corrigível
com a granularidade anual do painel; a winsorização no percentil 99 e o
tratamento da UTI em duas partes o mitigam, e ele deve ser lembrado na
leitura dos anos de 2020 e 2021. O cuidado adotado: analisar a
ocupação como razão positiva, tratar a UTI em duas partes (ter ou não
UTI ativa; e, entre os que têm, a intensidade de uso) e aparar a cauda
extrema (winsorização no percentil 99: 125\% na internação e 171\% na
UTI), com atenção redobrada ao biênio pandêmico.\footnote{Duas
alternativas metodológicas ficam registradas para trabalhos futuros: a
sensibilidade da winsorização a cortes no percentil 95 e no percentil
98, e a modelagem estrutural das diárias como contagem (Binomial
Negativa) com o logaritmo dos leitos como termo de exposição, que
dispensa a razão de ocupação e o seu denominador problemático.}

\textbf{Denominadores frágeis na pandemia.} Ao excluir as internações
COVID, alguns hospitais poderiam ficar com pouquíssimas internações
restantes, tornando suas taxas instáveis. No painel isso ocorre em
\textbf{apenas 2 observações, ambas do Hospital Infante D.\ Henrique
(CNES 2097613, São José do Rio Preto)}: 13 saídas não COVID em 2020 e
somente 4 em 2021, ano em que o TMP calculado degenera para zero (o
único zero de TMP de todo o painel). As duas observações recebem
marcação de denominador frágil e ficam de fora das análises de
resultado.

\textbf{Ruído estatístico dos hospitais pequenos.} A
Figura~\ref{fig:funil} organiza essa discussão: cada ponto é um
hospital num ano, posicionado pelo volume de saídas (eixo horizontal) e
pela mortalidade (eixo vertical); as linhas mostram a faixa em que a
mortalidade deveria variar se as diferenças fossem puro acaso
estatístico, faixa que é larga para hospitais pequenos e estreita para
os grandes. O resultado: \textbf{cerca de 66\% das observações caem
fora até da faixa mais generosa} (1.033 acima e 1.051 abaixo). Isso diz
duas coisas ao mesmo tempo: as diferenças de mortalidade entre
hospitais são reais e estruturais, ligadas ao perfil assistencial, e
não flutuação de sorte; e, abaixo de mil saídas por ano, o acaso passa
a pesar de verdade, o que justifica a cautela extra com hospitais
pequenos.

\textbf{Valores extremos e completude.} Pela regra usual de detecção
(3 intervalos interquartílicos acima do terceiro quartil), há 18
observações extremas de mortalidade (acima de 17,2\%), 87 de faturamento
(acima de R\$~4.223 em valores correntes, ver
Seção~\ref{sec:tendencia}), 178 de alta complexidade e 33 de TMP. Os
extremos de permanência são poucos porque a subpopulação de longa
permanência que os concentraria saiu na exclusão dos hospitais
estruturalmente incomparáveis. Quanto à
completude, \textbf{o painel não tem nenhum valor faltante nos sete
indicadores} nem nas variáveis de estratificação; as únicas lacunas
estão em campos cadastrais (nome fantasia e especialização), sem
consequência para as análises.

\comoler{o funil é estreito à direita (hospitais grandes, onde o acaso
pesa pouco) e largo à esquerda (hospitais pequenos, onde uma morte a
mais ou a menos muda a taxa). Pontos fora das linhas curvas indicam
mortalidade que o acaso não explica: a nuvem inteira transbordando das
linhas mostra que as diferenças entre hospitais são estruturais, e a
região à esquerda de mil saídas mostra onde as taxas devem ser lidas
com reserva.}

\begin{figure}[H]
\centering
\includegraphics[width=.62\textwidth]{fig_ae_16_funnel_mortalidade.png}
\caption{Funil da mortalidade: cada ponto é um hospital num ano; a
linha horizontal é a mortalidade média da rede (5,57\%) e as linhas
curvas delimitam a faixa esperada por puro acaso estatístico (95\% e
99,8\%), conforme o volume de saídas.}
\label{fig:funil}
\end{figure}

\comofeito{gráfico de dispersão (funil). Cada ponto é um hospital em um
ano, posicionado pelo \textbf{volume de saídas} no eixo horizontal (em
escala logarítmica) e pela \textbf{mortalidade geral} no vertical. A
linha horizontal é a mortalidade média da rede (5,57\%) e as duas linhas
curvas são os limites dentro dos quais a mortalidade deveria variar se as
diferenças fossem só acaso estatístico (95\% e 99,8\%), calculados a
partir do volume de cada hospital, por isso o funil é largo à
esquerda (hospitais pequenos) e estreito à direita (grandes).}
